\chapter{Analiza domeniului și studiu comparativ}

\section{Evoluția spațiului design-to-code}

Zona de \textit{design-to-code} nu este nouă, dar a devenit mult mai vizibilă odată cu maturizarea aplicațiilor web și cu apariția modelelor generative moderne. În literatura mai veche, problema era formulată deseori ca transformarea unei descrieri vizuale într-un cod de interfață. Un exemplu cunoscut este \textit{pix2code}, unde autorul tratează generarea codului din capturi de ecran ale interfeței \cite{beltramelli2017pix2code}. Mai recent, lucrările dedicate modelelor mari de limbaj și modelelor multimodale arată că interesul s-a mutat de la simpla reconstrucție a aspectului spre păstrarea structurii, a layout-ului și a intenției de produs \cite{shen2024uicoder,xu2025latcoder}.

În practică, piața a evoluat într-o direcție puțin diferită de lucrările academice. Produsele comerciale au observat că utilizatorul nu are nevoie doar de un model care produce cod, ci de un întreg flux de lucru: proiectare, organizare pe pagini, colaborare, verificare, export, publicare și iterație. Din acest motiv, instrumentele moderne nu mai pot fi analizate doar prin întrebarea ``generează sau nu cod?'', ci prin ansamblul de funcții pe care îl oferă.

Mai există și un detaliu important. În multe situații, problema nu este că nu se poate genera ceva, ci că rezultatul generat nu este suficient de ușor de integrat într-un produs real. Un layout poate arăta bine ca demonstrație, dar dacă nu are o structură logică, dacă nu respectă constrângeri tehnice și dacă nu poate fi controlat prin iterații mici, valoarea lui practică scade. Această observație a influențat direct proiectarea Favigon: în loc să trateze AI-ul ca pe un generator final, platforma îl include într-un lanț controlat, în care reprezentarea intermediară și validarea joacă un rol central.

\section{Criterii de comparație}

Pentru a poziționa corect proiectul, au fost alese criterii relevante atât din punct de vedere tehnic, cât și din punct de vedere al utilității pentru utilizator:

\begin{itemize}[leftmargin=1.5cm]
  \item \textbf{editare vizuală directă}: existența unui spațiu de lucru de tip canvas sau editor vizual;
  \item \textbf{asistență AI}: posibilitatea de a porni de la un prompt, sugestii sau generare automată de structură;
  \item \textbf{design-to-code}: existența unei conversii către cod sau măcar a unor artefacte pentru handoff tehnic;
  \item \textbf{orientare spre produs final}: suport pentru publicare, organizare pe pagini sau dezvoltare continuă;
  \item \textbf{partajare și colaborare}: mecanisme prin care rezultatele pot fi distribuite sau reutilizate;
  \item \textbf{control tehnic al ieșirii}: cât de mult poate fi influențată structura rezultatului și cât de ușor poate fi integrat în cod existent.
\end{itemize}

Aceste criterii nu sunt perfecte și nu încearcă să stabilească ``cel mai bun'' produs. Scopul lor este mai modest: să arate unde se află Favigon pe o hartă a instrumentelor existente și ce gol încearcă să acopere.

\section{Aplicații relevante din aceeași zonă}

\subsection{Figma}

Figma este, în mod firesc, un punct de referință pentru orice discuție despre proiectarea interfețelor. Din perspectiva acestei lucrări, interesul nu este dat doar de partea de design colaborativ, ci și de Dev Mode. Documentația oficială arată clar că Dev Mode este gândit ca spațiu pentru inspectarea designurilor și pentru apropierea dintre design și implementare \cite{figmaDevMode2026}. Cu toate acestea, Figma rămâne în primul rând un instrument de design și handoff. Chiar dacă poate genera fragmente de cod sau poate conecta design system-uri la bazele de cod, el nu devine automat un mediu complet de construire și publicare a produsului.

Pentru Favigon, Figma este util mai ales ca termen de comparație în două direcții. Prima direcție este ideea de structură vizuală clară și multi-page. A doua este faptul că handoff-ul către dezvoltare nu este suficient în sine; dacă se urmăresc iterații rapide, este necesar un spațiu în care designul să poată fi nu doar inspectat, ci și transformat mai departe în ceva executabil.

\subsection{Webflow}

Webflow este relevant deoarece reprezintă foarte bine zona de \textit{visual website builder}. Pagina oficială pune accent pe construcția vizuală, publicare și generare rapidă de site-uri \cite{webflow2026}. Avantajul evident al acestei abordări este că rezultatul ajunge repede online și este ușor de gestionat din punct de vedere al conținutului. În schimb, Webflow are un focus clar pe zona de website publishing și CMS, nu pe un model general de interfață reutilizabil între mai multe framework-uri.

Pentru poziționarea Favigon, Webflow este important tocmai pentru a delimita intenția aplicației. Nu a fost urmărită construirea încă unui publicator vizual de site-uri. Accentul cade mai mult pe controlul tehnic al structurii, exportul în cod și legătura dintre un editor vizual și o reprezentare internă care poate fi convertită.

\subsection{Framer}

Framer este interesant fiindcă mută discuția mai aproape de zona AI. Pagina oficială pentru Wireframer arată o abordare în care un prompt este folosit pentru a construi rapid un layout responsive, cu accent pe structură și flux \cite{framerWireframer2026}. Este o direcție apropiată de cea urmărită și aici, însă diferența majoră este că în Favigon este făcută explicită separarea dintre intenție, structură și stilizare.

Cu alte cuvinte, Framer demonstrează că generarea ghidată de prompt poate fi utilă pentru pornirea rapidă a unui proiect. Totuși, pentru o lucrare inginerească este mai valoros să fie explicați și controlați pașii interni ai generării, nu doar să fie observat rezultatul final.

\subsection{Builder.io}

Builder.io este probabil cel mai apropiat reper comercial pentru direcția urmărită în această lucrare. Platforma lor pune accent pe dezvoltare vizuală, AI și reutilizarea componentelor existente \cite{builderPlatform2026,builderAi2026}. Ideea de a nu porni de la zero de fiecare dată, ci de a avea o punte între design, componente și cod, este una foarte relevantă.

Diferența principală este că Builder.io este orientat puternic spre integrarea într-un ecosistem existent și spre productivitate în contexte enterprise. Favigon, în schimb, este un proiect academic și experimental, în care accentul este pus pe claritatea modelului intern și pe trasabilitatea tehnică a fiecărei etape. Asta înseamnă că produsul nu concurează direct cu Builder.io la nivel de maturitate comercială, dar tratează o parte din probleme într-un mod mai transparent pentru analiză și explicație.

\section{Tabel comparativ}

\begin{table}[H]
  \centering
  \caption{Comparație între Favigon și aplicații relevante}
  \label{tab:comparatie-platforme}
  \renewcommand{\arraystretch}{1.2}
  \begin{tabularx}{\textwidth}{>{\RaggedRight\arraybackslash}p{3.1cm}YYYYY}
    \toprule
    \textbf{Criteriu} & \textbf{Figma} & \textbf{Webflow} & \textbf{Framer} & \textbf{Builder.io} & \textbf{Favigon} \\
    \midrule
    Editor vizual direct & Da & Da & Da & Da & Da \\
    Pornire din prompt AI & Parțial & Da & Da & Da & Da \\
    Model intern orientat spre export multi-framework & Nu explicit & Nu & Limitat & Parțial & Da \\
    Export către HTML/React/Angular & Nu ca obiectiv principal & Nu în forma urmărită aici & Limitat & Orientat pe integrare & Da \\
    Funcții sociale în aceeași platformă & Nu & Nu & Nu & Nu & Da \\
    Publicare proiecte proprii & Parțial & Da & Da & Da & Parțial \\
    Reutilizare prin fork / star / like & Nu & Nu & Nu & Nu & Da \\
    Accent pe controlul etapelor AI & Nu & Nu & Parțial & Parțial & Da \\
    \bottomrule
  \end{tabularx}
\end{table}

Tabelul~\ref{tab:comparatie-platforme} nu are rolul de a demonstra că Favigon este ``mai bun'' decât produsele consacrate. Ar fi o comparație nerealistă. Scopul lui este să arate că Favigon combină, într-o formă academică și experimentală, elemente care în alte produse apar separate: editor vizual, asistență AI, export în cod și funcții sociale.

\section{Poziționarea proiectului Favigon}

În urma comparației, Favigon poate fi poziționat ca o platformă experimentală de tip \textit{AI-assisted visual builder}, orientată mai puțin spre publicare imediată și mai mult spre explorarea unui flux coerent între idee, design, model intern și cod executabil. Cu alte cuvinte, produsul nu își propune să concureze prin maturitatea ecosistemului, ci prin claritatea traseului tehnic pe care îl oferă.

În termeni practici, produsul este mai potrivit pentru prototipare rapidă, experimente de interfață, proiecte academice și aplicații mici sau medii decât pentru scenarii enterprise foarte rigide. Această delimitare este importantă, deoarece explică mai bine valoarea reală a platformei: nu promite înlocuirea unui ecosistem comercial matur, ci oferă un cadru integrat pentru a trece mai repede de la idee la un rezultat executabil și analizabil tehnic.

Această poziționare vine și cu limite. De exemplu, Favigon nu are nivelul de rafinament vizual sau de integrare enterprise al unor produse comerciale mature. Nici partea de colaborare simultană în timp real nu este încă implementată. Totuși, pentru scopul lucrării, aceste limite nu sunt un dezavantaj major. Ele permit concentrarea pe întrebările urmărite aici:

\begin{itemize}[leftmargin=1.5cm]
  \item cum pot reprezenta o interfață într-o formă suficient de stabilă încât să o editez și să o export;
  \item cum pot integra AI într-un mod controlabil;
  \item cum pot susține reutilizarea prin funcții sociale simple;
  \item cum pot organiza întregul produs astfel încât să rămână clar tehnic.
\end{itemize}

\section{Golul pe care îl acoperă Favigon}

Din analiza de mai sus rezultă că există un spațiu intermediar între instrumentele clasice de design, generatoarele AI rapide și platformele orientate strict spre publicare. Favigon încearcă să acopere tocmai acest spațiu intermediar. Mai exact, platforma propune:

\begin{enumerate}[leftmargin=1.5cm]
  \item un editor propriu, nu doar o zonă de import sau inspectare;
  \item un model intermediar explicit, nu doar un rezultat final opac;
  \item export multi-framework, nu doar un cod specific unei singure ținte;
  \item funcții sociale integrate, astfel încât proiectele să poată circula în comunitate;
  \item o integrare a AI-ului care poate fi analizată, validată și îmbunătățită pas cu pas.
\end{enumerate}

Din punct de vedere academic, acesta reprezintă un unghi bun pentru o lucrare de licență. Proiectul nu se reduce la folosirea unui model AI pentru obținerea unor pagini, ci pune în discuție mai multe aspecte inginerești: modelare internă, arhitectură pe straturi, validare, testare, UX și decizii de produs.

\section{Concluzii intermediare}

Analiza domeniului arată că proiectul Favigon nu apare într-un vid, ci într-un context în care instrumentele pentru design, dezvoltare vizuală și AI cresc rapid. În același timp, ea arată și de ce proiectul are sens ca temă de licență. Există deja instrumente puternice, dar nu toate tratează împreună editarea vizuală, generarea controlată, exportul tehnic și partajarea în aceeași platformă.

Concluzia acestui capitol este că Favigon nu trebuie evaluat ca un produs comercial finalizat, ci ca o platformă software care încearcă să demonstreze fezabilitatea unui flux integrat. Pe baza acestei concluzii, capitolul următor trece de la piață și comparații la fundamentele teoretice și tehnologiile pe care s-a sprijinit implementarea.
